\chapter{Construcción de espacios topológicos}
\section{Topología inicial}
Dado un espacio topológico $\displaystyle \left(X, \mathcal{T}\right) $ y una función $\displaystyle f : Y \to \left(X, \mathcal{T}\right) $, nos preguntamos si hay alguna topología $\displaystyle \mathcal{T}' $ en $\displaystyle Y $ que haga que $\displaystyle f $ sea continua. Es trivial que si tomamos $\displaystyle \mathcal{T}' = \mathcal{P}\left(Y\right) $ obtenemos el resultado deseado. En general, buscaremos la topología menos fina que cumpla lo requerido. 
 
\begin{observation}
Si $\displaystyle f $ es continua tendremos que para todo abierto $\displaystyle U \in \mathcal{T} $ debe ser que $\displaystyle f^{-1}\left(U\right)\in \mathcal{T}' $. Por tanto, la topología que buscamos es $\displaystyle \left\{ f^{-1}\left(U\right) \; : \; U \in \mathcal{T}\right\}  $, que es una topología en $\displaystyle Y $.
\end{observation}
\begin{proof}
	Veamos que la familia $\displaystyle \mathcal{T}^{f}=\left\{ f^{-1}\left(U\right) \; : \; U \in \mathcal{T}\right\}  $ es una topología. 
	\begin{enumerate}
	\item Tenemos que $\displaystyle \emptyset = f^{-1}\left(\emptyset \right) $ y $\displaystyle Y = f^{-1}\left(X\right) $, por lo que $\displaystyle \emptyset  $ y $\displaystyle X $ están en $\displaystyle \mathcal{T}^{f} $.  
	\item Supongamos que $\displaystyle \left\{ U_{i}\right\} _{i \in I}\subset \mathcal{T}^{f}$, entonces existe $\displaystyle V_{i} $ tal que $\displaystyle U_{i} = f^{-1}\left(V_{i}\right) $ para cada $\displaystyle i \in I $. Como $\displaystyle \bigcup_{i \in I}V_{i} \in \mathcal{T} $ tendremos que  
		\[\bigcup_{i \in I}U_{i}=\bigcup_{i \in I}f^{-1}\left(V_{i}\right)=f^{-1}\left(\bigcup_{i \in I}V_{i}\right) \in \mathcal{T}^{f} .\]
	\item Supongamos que $\displaystyle U_{1}, U_{2} \in \mathcal{T}^{f} $. Entonces existen $\displaystyle V_{1}, V_{2} \in \mathcal{T} $ tales que $\displaystyle U_{i} = f^{-1}\left(V_{i}\right) $ para $\displaystyle i \in \left\{ 1,2\right\}  $. De esta forma, 
		\[U_{1} \cap U_{2}=f^{-1}\left(V_{1}\right)\cap f^{-1}\left(V_{2}\right)=f^{-1}\left(V_{1} \cap V_{2}\right)\in \mathcal{T}^{f} .\]
	\end{enumerate}
	Veamos que $\displaystyle \mathcal{T}^{f} $ es la topología más fina que cumple lo que queremos. Si tenemos otra topología $\displaystyle \mathcal{T}' $ en $\displaystyle Y $ tal que $\displaystyle f : \left(Y, \mathcal{T}'\right)\to \left(X, \mathcal{T}\right) $ es continua, tendremos que para cada $\displaystyle U \in \mathcal{T} $, $\displaystyle f^{-1}\left(U\right)\in \mathcal{T}' $, de forma que $\displaystyle \mathcal{T}^{f}\subset \mathcal{T}' $. 
\end{proof}

\begin{definition}[Topología imagen inversa]
Sea $\displaystyle f : X \to Y $ una aplicación y consideremos fijada en $\displaystyle X $ una topología $\displaystyle \mathcal{T} $. La familia de subconjuntos
\[ \left\{ f^{-1}\left(U\right) \; : \; U \in \mathcal{T}\right\}  \]
define una topología que llamaremos \textbf{topología imagen inversa por $\displaystyle f $}, o \textbf{topología inicial para $\displaystyle f $} y la denotaremos por $\displaystyle f^{-1}\left(\mathcal{T}\right) $ o $\displaystyle \mathcal{T}^{f} $. 
\end{definition}
\begin{prop}[Propiedad universal]
Sea $\displaystyle f : Y \to X $ una aplicación y $\displaystyle \mathcal{T} $ una topología fijada en $\displaystyle X $. Las siguientes afirmaciones son equivalentes:
\begin{enumerate}
\item La topología $\displaystyle \mathcal{T}' $ en $\displaystyle Y $ es la topología inicial para $\displaystyle f $. 
\item Para cada espacio topológico $\displaystyle \left(Z, \mathcal{T}''\right) $ y para cada aplicación $\displaystyle g : \left(Z, \mathcal{T}''\right)\to \left(Y, \mathcal{T}'\right) $ se tiene que $\displaystyle g $ es continua si y sólo si 
	\[f\circ g : \left(Z, \mathcal{T}''\right)\to \left(X, \mathcal{T}\right) \]
	es continua.
\end{enumerate}
\end{prop}
\begin{proof}
Demostramos cada implicación.
\begin{description}
\item[(1) $\displaystyle \Rightarrow $ (2)] Supongamos que $\displaystyle \mathcal{T}' = \mathcal{T}^{f} $. Entonces, dado un espacio topológico $\displaystyle \left(Z, \mathcal{T}''\right) $ y una aplicación
	\[g : \left(Z, \mathcal{T}''\right)\to \left(Y, \mathcal{T}'\right) ,\]
	se tiene:
	\begin{itemize}
	\item Si $\displaystyle g $ es continua, dado un abierto $\displaystyle U \subset X $, tendremos que 
		\[\left(f\circ g\right)^{-1}\left(U\right)=g^{-1}\left(f^{-1}\left(U\right)\right)\in \mathcal{T}'' ,\]
		puesto que $\displaystyle f^{-1}\left(U\right)\in \mathcal{T}' =\mathcal{T}^{f} $. Por tanto, $\displaystyle f\circ g $ es continua. 
	\item Si $\displaystyle f\circ g $ es continua, entonces para cada $\displaystyle V \in \mathcal{T}'=\mathcal{T}^{f} $, se verifica que existe $\displaystyle U \in \mathcal{T} $ tal que $\displaystyle V = f^{-1}\left(U\right) $. Así, 
		\[g^{-1}\left(V\right)=g^{-1}\left(f^{-1}\left(U\right)\right)=\left(f\circ g\right)^{-1}\left(U\right)\in \mathcal{T}'' ,\]
		por lo que $\displaystyle g $ es continua. 
	\end{itemize}
\item[(2) $\displaystyle \Rightarrow $ (1)] Tomando como $\displaystyle \left(Z, \mathcal{T}''\right) $ el espacio $\displaystyle \left(Y, \mathcal{T}^{f}\right) $ y tomando $\displaystyle g $ como la aplicación identidad en $\displaystyle Y $ tendremos que 
	\[f\circ g = f\circ id_Y = f : \left(Y, \mathcal{T}^{f}\right)\to \left(X, \mathcal{T}\right) .\]
	Claramente esta función es continua. Pero entonces, 
	\[g = id_Y : \left(Y, \mathcal{T}^{f}\right)\to \left(Y, \mathcal{T}'\right) ,\]
	es continua por lo que debe ser que $\displaystyle \mathcal{T}' \subset \mathcal{T}^{f} $. Por otro lado, tomando como $\displaystyle \left(Z, \mathcal{T}''\right) $ el espacio $\displaystyle \left(Y, \mathcal{T}'\right) $ y como $\displaystyle g $ la identidad en $\displaystyle Y $, 
	\[g =id_Y : \left(Y, \mathcal{T}'\right) \to \left(Y, \mathcal{T}'\right) ,\]
	es claramente continua y por hipótesis sabemos que 
	\[f\circ g = f\circ id_Y = f : \left(Y, \mathcal{T}'\right)\to \left(X, \mathcal{T}\right) ,\]
	es continua por lo que $\displaystyle \mathcal{T}^{f}\subset \mathcal{T}' $. 
\end{description}
\end{proof}

\begin{eg}
Consideremos un espacio topológico $\displaystyle \left(X, \mathcal{T}\right) $ e $\displaystyle Y \subset X $. Con la aplicación $\displaystyle j : Y \to X : y \to y $, tendremos que 
\[\mathcal{T}^{j}= \left\{ j^{-1}\left(U\right) \; : \; U \in \mathcal{T}\right\} = \left\{ U \cap Y \; : \; U \in \mathcal{T}\right\}  ,\]
por lo que obtenemos la topología relativa. 
\end{eg}
\subsection{Embebimientos topológicos}

\begin{definition}[Embebimiento topológico]
Un \textbf{embebimiento topológico} es una aplicación inyectiva $\displaystyle f : \left(Y, \mathcal{T}'\right)\to \left(X, \mathcal{T}\right) $ tal que $\displaystyle \mathcal{T}' = \mathcal{T}^{f} $. 
\end{definition}
\begin{observation}
Si $\displaystyle f $ es embebimiento entonces es continua.
\end{observation}

\begin{eg}
El ejemplo anterior de la topología relativa es un embebimiento topológico. 
\end{eg}

\begin{prop}
Las siguientes afirmaciones son equivalentes. 
\begin{enumerate}
\item $\displaystyle f : \left(Y, \mathcal{T}'\right)\to \left(X, \mathcal{T}\right) $ es embebimiento topológico. 
\item $\displaystyle f : \left(Y, \mathcal{T}'\right)\to \left(f\left(Y\right), \mathcal{T}|_{f\left(Y\right)}\right) $ es homeomorfismo. 
\end{enumerate}
\end{prop}
\begin{proof}
Demostramos ambas implicaciones.
\begin{description}
\item[(1) $\displaystyle \Rightarrow $ (2)] Claramente tendremos que $\displaystyle f : Y \to f\left(Y\right) $ es biyectiva. Sea $\displaystyle U \in \mathcal{T} $ de forma que $\displaystyle U \cap f\left(Y\right)\in \mathcal{T}|_{f\left(Y\right)} $, entonces
	\[f^{-1}\left(U \cap f\left(Y\right)\right)= f^{-1}\left(U\right) \cap f^{-1}\left(f\left(Y\right)\right)=f^{-1}\left(U\right)\in \mathcal{T}' .\]
	Así, tenemos que $\displaystyle f $ es continua. Ahora, si cogemos $\displaystyle V \in \mathcal{T}' $ tendremos que existe $\displaystyle U \in \mathcal{T} $ tal que $\displaystyle V = f^{-1}\left(U\right) $. De esta forma, $\displaystyle f\left(V\right)=U \cap f\left(Y\right)\in \mathcal{T}|_{f\left(Y\right)} $ y tenemos que $\displaystyle f : Y \to f\left(Y\right) $ es abierta. Así, es homeomorfismo.
\item[(2) $\displaystyle \Rightarrow $ (1)] Si $\displaystyle f : \left(Y, \mathcal{T}'\right)\to \left(f\left(Y\right), \mathcal{T}|_{f\left(Y\right)}\right) $ es homeomorfismo tenemos que es biyectiva, por lo que $\displaystyle f:\left(Y, \mathcal{T}'\right)\to \left(X, \mathcal{T}\right) $ es inyectiva. Queremos ver que $\displaystyle \mathcal{T}' = \mathcal{T}^{f} $. Por un lado, dado $\displaystyle U \in \mathcal{T} $ tendremos que $\displaystyle U \cap f\left(Y\right)\in \mathcal{T}|_{f\left(Y\right)} $ por lo que 
	\[f^{-1}\left(U\cap f\left(Y\right)\right) = f^{-1}\left(U\right) \cap f^{-1}\left(f\left(Y\right)\right)=f^{-1}\left(U\right) \cap Y=f^{-1}\left(U\right) \in \mathcal{T}' .\]
	Así, debe ser que $\displaystyle \mathcal{T}^{f}\subset \mathcal{T}' $. Por otro lado, dado $\displaystyle U \in \mathcal{T}' $ tenemos que existe $\displaystyle V \in \mathcal{T} $ tal que $\displaystyle f\left(U\right)=V \cap f\left(Y\right) $, por lo que 
	\[U = f^{-1}\left(V \cap f\left(Y\right)\right)=f^{-1}\left(V\right)\cap f^{-1}\left(f\left(Y\right)\right)=f^{-1}\left(V\right) .\]
	Así, $\displaystyle \mathcal{T}' \subset \mathcal{T}^{f} $. 
\end{description}
\end{proof}

\begin{eg}
Consideremos $\displaystyle f : \left(\R, \mathcal{T}_{[,)}\right)\to \left(\R^{2}, \mathcal{T}\right): x \to \left(x,x\right) $ con $\displaystyle \mathcal{T} $ la topología del ejercicio 15 de la hoja 1. Esto es un embebimiento de la recta de Sofrengein en el espacio $\displaystyle \left(\R^{2}, \mathcal{T}\right) $.
En efecto, es sencillo ver que se trata de una función inyectiva y que $\displaystyle \mathcal{T}^{f}=\mathcal{T}_{[,)} $. 
\end{eg}
\begin{prop}[Condiciones suficientes para embebimiento topológico]
Dada una aplicación $\displaystyle f : \left(Y, \mathcal{T}'\right)\to \left(X, \mathcal{T}\right) $, se tiene:
\begin{enumerate}
\item Si $\displaystyle f $ es inyectiva, continua y abierta, entonces $\displaystyle f $ es embebimiento topológico. 
\item Si $\displaystyle f $ es inyectiva, continua y cerrada, entonces $\displaystyle f $ es embebimiento topológico.
\end{enumerate}
\end{prop}
\begin{proof}
Supongamos que $\displaystyle f : \left(Y, \mathcal{T}'\right)\to \left(X, \mathcal{T}\right) $ es una aplicación entre espacios topológicos. Demostramos sólo \textbf{(1)} puesto que \textbf{(2)} es análogo. Tenemos que para cada $\displaystyle V \in \mathcal{T}' $, $\displaystyle f\left(V\right)\in \mathcal{T} $. Como $\displaystyle f $ es abierta tendremos que $\displaystyle f\left(V\right)\subset f\left(Y\right) $ es abierto en $\displaystyle f\left(Y\right) $. De esta forma, la función $\displaystyle f : \left(Y, \mathcal{T}'\right)\to \left(f\left(Y\right), \mathcal{T}|_{f\left(Y\right)}\right) $ es biyectiva, continua y abierta, por lo que es homeomorfismo. Por la proposición anterior tendremos que $\displaystyle f $ es embebimiento topológico.
\end{proof}
\begin{observation}
El recíproco en general no es cierto, basta con ver el siguiente ejemplo.
\end{observation}

\begin{eg}
	Consideremos una aplicación $\displaystyle f : \left(\left[0,1\right] , \mathcal{T}_{u}|_{\left[0,1\right] }\right)\to\left(\left(-1,2\right), \mathcal{T}_{u}|_{\left(-1,2\right)}\right) $ que es la inyección $\displaystyle j $ que vimos antes. Es un embebimiento topológico puesto que es inyectiva y $\displaystyle \mathcal{T}^{f}= \mathcal{T}_{u}|_{\left[0,1\right] } $. Sin embargo no es una aplicación abierta puesto que $\displaystyle \left[0,1\right]  $ no es abierto en $\displaystyle \mathcal{T}_{u}|_{\left(-1,2\right)} $.
\end{eg}

\begin{prop}
Dado un embebimiento topológico $\displaystyle f : Y \to X $ se tiene que:
\begin{enumerate}
\item $\displaystyle f $ es abierta si y sólo si $\displaystyle f\left(Y\right) $ es abierto en $\displaystyle X $. 
\item $\displaystyle f $ es cerrada si y sólo si $\displaystyle f\left(Y\right) $ es cerrado en $\displaystyle X $. 
\end{enumerate}
\end{prop}
\begin{proof}
Demostramos sólo \textbf{(1)} puesto que la demostración de \textbf{(2)} es análoga. Como $\displaystyle f $ es embebimiento topológico tendremos que la topología en $\displaystyle Y $ es $\displaystyle \mathcal{T}^{f} $. 
Si $\displaystyle f $ es abierta tenemos que para todo $\displaystyle V \in \mathcal{T}^{f} $ se tiene que $\displaystyle f\left(V\right)\in \mathcal{T} $. Basta con tomar $\displaystyle V = Y \in \mathcal{T}^{f} $ (porque $\displaystyle Y =f^{-1}\left(X\right) $, $\displaystyle X \in \mathcal{T} $) para ver que $\displaystyle f\left(Y\right)\in \mathcal{T} $.
Recíprocamente, supongamos que $\displaystyle f\left(Y\right) $ es abierto en $\displaystyle X $. Como para cada abierto $\displaystyle V \subset Y $ se tiene que $\displaystyle V = f^{-1}\left(U\right) $ abierto en $\displaystyle X $, tendremos que $\displaystyle f\left(V\right)=U \cap f\left(Y\right) $ es abierto en $\displaystyle X $. 
\end{proof}

\begin{observation}
Los embebimientos topológicos permiten considerar algunos espacios topológicos como subespacios de otros. Que un espacio sea subespacio de otro quiere decir que existe un embebimiento topológico del primero al segundo.
\end{observation}

\subsection{Productos (finitos) de espacios topológicos}
Consideremos un número finito de espacios topológicos $\displaystyle \left\{ \left(X_{i}, \mathcal{T}_{i}\right)\right\}_{i = 1, \ldots, r} $ y el conjunto 
\[Y = \prod^{r}_{i = 1}X_{i} = \left\{ \left(x_{1}, \ldots, x_{r}\right) \; : \; x_{i} \in X_{i}, i \in \left\{ 1, \ldots, r\right\} \right\}  .\]
Consideramos también las proyecciones
\[p_{j} : \prod^{r}_{i = 1}X_{i} \to X_{j} : \left(x_{1}, \ldots, x_{r}\right) \to x_{j} \]
para cada $\displaystyle j \in \left\{ 1, \ldots, r\right\}  $. Buscamos la topología menos fina en $\displaystyle \prod^{r}_{i = 1}X_{i} $ que haga que todas las proyecciones sean continuas. Una posible solución sería coger el conjunto
\[ \mathcal{S} = \left\{ p_{j}^{-1}\left(U_{j}\right) \; : \; j \in \left\{ 1, \ldots, r\right\} , U_{j}\in \mathcal{T}_{j}\right\}  .\]
El problema es que este conjunto no es una topología. En efecto, tendremos que 
\[p_{j}^{-1}\left(U_{j}\right)=X_{1} \times \cdots \times X_{j-1} \times U_{j} \times X_{j + 1} \times \cdots \times X_{r} .\]
Se puede ver que la intersección de dos preimágenes $\displaystyle p_{1}^{-1}\left(U_{1}\right) \cap p^{-1}\left(U_{2}\right) $ no está en $\displaystyle \mathcal{S} $. 
	\begin{center}
    \includegraphics[width=0.4\textwidth]{~/Desktop/Images/Topologia4.png}
\end{center}
La solución consiste en coger una base 
\[\mathcal{B} = \left\{ U_{1} \times \cdots \times U_{r} \; : \; U_{i} \in \mathcal{T}_{i}, i \in \left\{ 1, \ldots, r\right\} \right\}  ,\]
que es la base determinada por la subbase $\displaystyle \mathcal{S} $. En efecto, tenemos que 
\[U_{1} \times \cdots \times U_{r}=p^{-1}_{1}\left(U_{1}\right) \cap \cdots \cap p_{r}^{-1}\left(U_{r}\right) .\] 
Veamos que $\displaystyle \mathcal{B} $ es en efecto una base:
\begin{itemize}
\item Por un lado, como $\displaystyle X_{i}=\bigcup_{j \in J_{i}}U_{j} $ con $\displaystyle U_{j}\in \mathcal{T}_{i} $, tendremos que 
	\[\prod^{r}_{i = 1}X_{i}=\bigcup_{U_{j}^{i}\in \mathcal{T}_{i}, i \in \left\{ 1, \ldots, r\right\} , j \in J} U_{j}^{1} \times \cdots \times U_{j}^{r},.\]
\item Por otro lado, supongamos que $\displaystyle B_{1} = U_{1} \times \cdots \times U_{r}\in \mathcal{B} $ y $\displaystyle B_{2} = V_{1} \times \cdots \times V_{r} \in \mathcal{B} $ no son disjuntos, por lo que existe $\displaystyle \left(x_{1}, \ldots, x_{r}\right)\in B_{1} \cap B_{2} $. Tendremos que 
	\[\left(x_{1}, \ldots, x_{r}\right)\in \left(U_{1} \cap V_{1}\right) \times \cdots \times \left(U_{r} \cap V_{r}\right) \subset B_{1} \cap B_{2}  .\]
	Además, es sencillo ver que $\displaystyle \left(U_{1} \cap V_{1}\right) \times \cdots \times \left(U_{r} \cap V_{r}\right)\in \mathcal{B} $. 
\end{itemize}
Hemos visto que $\displaystyle \mathcal{B} $ es una base, ahora veamos que genera la topología menos fina que hace continuas todas las proyecciones $\displaystyle p_{j} $ para $\displaystyle j \in \left\{ 1, \ldots, r\right\}  $. Si $\displaystyle \mathcal{T} $ es una topología que cumple estas características, tendremos que si $\displaystyle U_{j} \in \mathcal{T}_{j} $ entonces
\[p_{j}^{-1}\left(U_{j}\right)= X_{1} \times \cdots \times X_{j-1} \times U_{j} \times X_{j+1} \times \cdots \times X_{r} \in \mathcal{T} .\]
Así, para $\displaystyle U = U_{1} \times \cdots \times U_{r} \in \mathcal{T}_{\mathcal{B}}$ tendremos que $\displaystyle U = p^{-1}_{1}\left(U_{1}\right) \cap \cdots \cap p^{-1}_{r}\left(U_{r}\right) \in \mathcal{T} $, por lo que $\displaystyle \mathcal{T}_{\mathcal{B}}\subset \mathcal{T} $.  
\begin{definition}[Topología producto]
	Dado un conjunto finito de espacios topológicos $\displaystyle \left\{ \left(X_{i}, \mathcal{T}_{i}\right)\right\} _{i = 1, \ldots, r} $, a la topología determinada por la subbase 
	\[\mathcal{S}= \left\{ p_{j}^{-1}\left(U_{j}\right) \; : \; j \in \left\{ 1, \ldots, r\right\} , U_{j}\in \mathcal{T}_{j}\right\} .\]
	la denominamos \textbf{topología producto} y la escribimos $\displaystyle \prod^{r}_{i = 1}\mathcal{T}_{i} $. 
\end{definition}
 
\begin{eg}
Consideremos el producto de $\displaystyle n $ copias de $\displaystyle \left(\R, \mathcal{T}_{u}\right) $. Es sencillo ver que la topología producto en $\displaystyle \R^{n} $ es la topología usual en $\displaystyle \R^{n} $. En efecto, podemos ver que $\displaystyle \left(a_{1}, b_{1}\right)\times \cdots \times\left(a_{n}, b_{n}\right) $ con $\displaystyle \left(a_{i}, b_{i}\right) $ intervalo de $\displaystyle \R $, es una base de la topología usual en $\displaystyle \R^{n} $. 
\end{eg}

\begin{prop}
Se consideran espacios topológicos $\displaystyle \left(X_{i}, \mathcal{T}_{i}\right) $ para $\displaystyle i = 1, \ldots, r $, el conjunto producto $\displaystyle Y = \prod^{r}_{i=1}X_{i} $ y las aplicaciones proyección $\displaystyle p_{i} : Y \to X_{i} $. Las siguientes afirmaciones son equivalentes:
\begin{enumerate}
\item La topología $\displaystyle \mathcal{T}' $ en $\displaystyle Y $ es la topología producto.
\item Para cada espacio topológico $\displaystyle \left(Z, \mathcal{T}''\right) $ y para cada aplicación 
	\[g = \left(g_{1}, \ldots, g_{r}\right) : \left(Z, \mathcal{T}''\right)\to \left(Y, \mathcal{T}'\right) ,\]
	se tiene que $\displaystyle g $ es continua si y sólo si todas las aplicaciones 
	\[g_{i} = p_{i}\circ g : \left(Z, \mathcal{T}''\right)\to \left(X_{i}, \mathcal{T}_{i}\right), \; i = 1, \ldots, r ,\]
	son continuas.
\end{enumerate}
\end{prop}
\begin{proof}
Consideramos ambas implicaciones.
\begin{description}
\item[($\displaystyle \Rightarrow $)] Supongamos que $\displaystyle \mathcal{T}' = \prod^{r}_{i = 1}\mathcal{T}_{i} $. 
	\begin{itemize}
		\item Si $\displaystyle g $ es continua tendremos que como $\displaystyle p_{j} $ es continua $\displaystyle \forall j \in \left\{ 1, \ldots, r\right\}  $, tendremos que $\displaystyle p_{j}\circ g $ es continua. 
		\item Si $\displaystyle p_{j}\circ g $ es continua para cada $\displaystyle j \in \left\{ 1, \ldots, r\right\}  $, cogemos $\displaystyle U_{1} \times \cdots \times U_{r} \in \mathcal{T}' $, entonces tendremos que 
			\[g^{-1}\left(U_{1} \times \cdots \times U_{r}\right)= \bigcap_{i = 1}^{r}g_{i}^{-1}\left(U_{i}\right) = \bigcap_{i = 1}^{r}\left(p_{i}\circ g\right)^{-1}\left(U_{i}\right) \in \mathcal{T}'' .\]
		Así, tenemos que $\displaystyle g $ es continua.
	\end{itemize}
\item[($\displaystyle \Leftarrow $)] Elegimos $\displaystyle \left(Z, \mathcal{T}''\right)=\left(\prod^{r}_{i = 1}X_{i}, \prod^{r}_{i = 1}\mathcal{T}_{i}\right) $ y $\displaystyle g = id $. Entonces, 
	\[p_{j}\circ id = p_{j} : \left(\prod^{r}_{i = 1}X_{i}, \prod^{r}_{i = 1}\mathcal{T}_{i}\right) \to \left(X_{j}, \mathcal{T}_{j}\right) \]
	es continua, por lo que debe ser que 
	\[g = id : \left(\prod^{r}_{i = 1}X_{i}, \prod^{r}_{i = 1}\mathcal{T}_{i}\right)\to \left(\prod^{r}_{i = 1}X_{i}, \mathcal{T}'\right) \]
	es continua y por tanto tenemos que $\displaystyle \mathcal{T}' \subset \prod^{r}_{i = 1}\mathcal{T}_{i} $. Recíprocamente, cogiendo $\displaystyle \left(Z, \mathcal{T}''\right)=\left(\prod^{r}_{i = 1}X_{i}, \mathcal{T}'\right) $, cogemos $\displaystyle g = id $ que claramente es continua. Por tanto, $\displaystyle p_{j} $ es continua para $\displaystyle j \in \left\{ 1, \ldots, r\right\}  $ y debe ser que $\displaystyle \prod^{r}_{i = 1}\mathcal{T}_{i} \subset \mathcal{T}' $ \footnote{Estamos usando el hecho de que $\displaystyle \prod^{r}_{i = 1}\mathcal{T}_{i} $ es la topología menos fina que hace que todas las proyecciones sean continuas.}. 
\end{description}
\end{proof}
\begin{prop}
	Para cada $\displaystyle j = 1, \ldots, r $ las aplicaciones 
	\[ p_{j} : \left(\prod^{r}_{i = 1}X_{i}, \prod^{r}_{i = 1}\mathcal{T}_{i}\right) \to \left(X_{j}, \mathcal{T}_{j}\right)\]
	son abiertas.  
\end{prop}
\begin{proof}
	Supongamos que $\displaystyle \Omega \in \prod^{r}_{i = 1}\mathcal{T}_{i} $, entonces exsisten $\displaystyle U_{i}^{\alpha }\in \mathcal{T}_{i} $ con $\displaystyle \Omega = \bigcup_{\alpha \in \Lambda }\left(U_{1}^{\alpha } \times \cdots \times U_{r}^{\alpha }\right) $ y por tanto
	\[p_{j}\left(\Omega \right)=\bigcup_{\alpha \in \Lambda}p_{j}\left(U_{1}^{\alpha }\times \cdots \times U_{r}^{\alpha }\right) = \bigcup_{\alpha \in \Lambda }U_{j}^{\alpha} \in \mathcal{T}_{j} .\]
\end{proof}
\begin{observation}
Las proyecciones no son cerradas en general. Por ejemplo, consideremos en $\displaystyle \left(\R^{2}, \mathcal{T}_{u}\right)=\left(\R, \mathcal{T}_{u}\right) \times \left(\R, \mathcal{T}_{u}\right) $ la hipérbola $\displaystyle H = \left\{ xy = 1\right\} \subset \R^{2} $, que es un conjunto cerrado. Tenemos que $\displaystyle p_{1}\left(H\right)= \R / \left\{ 0\right\}  $ que no es cerrado.
\end{observation}

\begin{prop}
	Para cada $\displaystyle j \in \left\{ 1, \ldots, r\right\}  $ y para cualesquiera $\displaystyle a_{k}\in X_{k} $ con $\displaystyle k \neq j $, fijados la aplicación
	\[\alpha_{j} : \left(X_{j}, \mathcal{T}_{j}\right)\to \left(\prod^{r}_{i = 1}X_{i}, \prod^{r}_{i = 1}\mathcal{T}_{i}\right): x_{j}\to \left(a_{1}, \ldots, a_{j-1}, x_{j}, a_{j+1}, \ldots, a_{r}\right) ,\]
	es un embebimiento topológico.
\end{prop}

\begin{proof}
Consideremos $\displaystyle \alpha_{j} : \left(X_{j}, \mathcal{T}_{j}\right)\to \left(\alpha_{j}\left(X_{j}\right), \left(\prod^{r}_{i = 1}\mathcal{T}_{i}\right)|_{\alpha_{j}\left(X_{j}\right)}\right) $. Queremos ver que es un homeomorfismo. Claramente es biyectiva. Veamos ahora que es continua. Si cogemos $\displaystyle U_{1} \times \cdots \times U_{r} \in \prod^{r}_{i=1}\mathcal{T}_{i} $, distinguimos dos casos:
\begin{itemize}
	\item Si $\displaystyle \left(U_{1} \times \cdots \times U_{r}\right) \cap \alpha_{j}\left(X_{j}\right)=\emptyset $ es porque debe existir $\displaystyle k \in \left\{ 1, \ldots, r\right\} \backslash \left\{ j\right\}  $ con $\displaystyle a_{k}\not\in U_{k} $. En este caso $\displaystyle \alpha _{j}^{-1}\left(\emptyset\right)=\emptyset \in \mathcal{T}_{j} $. 
	\item Si no es vacío tendremos que $\displaystyle  $ 
		\[\left(U_{1}, \times \cdots \times U_{r}\right) \cap \alpha_{j}\left(X_{j}\right) = \left\{ a_{1}\right\} \times \cdots \times\left\{ a_{j - 1}\right\} \times U_{j} \times \left\{ a_{j+1}\right\} \times \cdots \times \left\{ a_{r}\right\}   .\]
De esta forma,  
\[\alpha ^{-1}_{j}\left(\left\{ a_{1}\right\} \times \cdots \times\left\{ a_{j - 1}\right\} \times U_{j} \times \left\{ a_{j+1}\right\} \times \cdots \times \left\{ a_{r}\right\}\right) = U_{j}\in \mathcal{T}_{j} .\]
\end{itemize}
Por tanto, $\displaystyle \alpha_{j} $ es continua. Finalmente, veamos que es abierta. Si $\displaystyle \Omega_{j}\in \mathcal{T}_{j} $ tendremos que 
\[
\begin{split}
	\alpha_{j}\left(\Omega_{j}\right)= & \left\{ a_{1}\right\}  \times \cdots \times \left\{ a_{j-1}\right\}  \times \Omega_{j} \times \left\{ a_{j+1}\right\} \times \cdots \times \left\{ a_{r}\right\} \\
	= &  \left(X_{1} \times \cdots \times X_{j-1} \times \Omega_{j} \times X_{j+1} \times \cdots \times X_{r}\right)\cap \alpha_{j}\left(X_{j}\right) \in  \left(\prod^{r}_{i = 1}X_{i}, \left(\prod^{r}_{i = 1}\mathcal{T}_{i}\right)|_{\alpha_{j}\left(X_{j}\right)}\right).
\end{split}
\]
Así, tenemos que $\displaystyle \alpha_{j} $ es abierta por lo que es homeomorfismo.
\end{proof}
\begin{eg}
Nos preguntamos si existen $\displaystyle \mathcal{T}_{1} $ y $\displaystyle \mathcal{T}_{2} $ topologías en $\displaystyle \R $ tales que $\displaystyle \left(\R^{2}, \mathcal{T}\right) = \left(\R, \mathcal{T}_{1}\right) \times \left(\R, \mathcal{T}_{2}\right) $ con $\displaystyle \mathcal{T} $ la topología del ejercicio 15. Necesariamente debe ser que $\displaystyle \mathcal{T}_{1} $ y $\displaystyle \mathcal{T}_{2} $ son la topología discreta. 
Sin embargo, esto no puede ser porque $\displaystyle \left(\R, \mathcal{T}_{\text{dis}}\right)\times\left(\R, \mathcal{T}_{\text{dis}}\right)=\left(\R^{2}, \mathcal{T}_{\text{dis}} \times \mathcal{T}_{\text{dis}}\right)\neq \left(\R^{2}, \mathcal{T}\right) $. 
\end{eg}
\begin{prop}
	Consideremos una familia de espacios topológicos $\displaystyle \left\{ X_{i}, \mathcal{T}_{i}\right\} _{i = 1, \ldots, r} $ y el espacio topológico producto $\displaystyle \left(\prod^{r}_{i = 1}X_{i}, \prod^{r}_{i = 1}\mathcal{T}_{i}\right) $. 
	\begin{enumerate}
		\item Si $\displaystyle V^{a_{j}}\subset X_{j} $ es entorno de $\displaystyle a_{j}\in X_{j} $ para $\displaystyle j \in \left\{ 1, \ldots, r\right\}  $, entonces
			\[V^{a_{1}}\times \cdots \times V^{a_{r}}\subset \prod^{r}_{i= 1}X_{i} \]
			es entorno del punto $\displaystyle \left(a_{1}, \ldots, a_{r}\right) $. 
		\item Si $\displaystyle \mathcal{V}\left(a_{i}\right) $ es una base de entornos de $\displaystyle a_{i}\in X_{i} $,
			\[\mathcal{V}\left(a_{1}, \ldots, a_{r}\right)= \left\{ V^{a_{1}} \times \cdots \times V^{a_{r}} \; : \; V^{a_{i}}\in \mathcal{V}\left(a_{i}\right), i = 1, \ldots, r\right\}  ,\]
			es una base de entornos de $\displaystyle \left(a_{1}, \ldots, a_{r}\right) $. 
		\item Si $\displaystyle \mathcal{B}_{i}\subset \mathcal{P}\left(X_{i}\right) $ es una base de $\displaystyle \mathcal{T}_{i} $ para $\displaystyle i = 1, \ldots, r $, entonces
			\[\mathcal{B}_{1 \times \cdots \times r} = \left\{ U_{1} \times \cdots \times U_{r} \; : \; U_{i} \in \mathcal{B}_{i}, i = 1, \ldots, r\right\}  ,\]
			es una base de $\displaystyle \prod^{r}_{i = 1}\mathcal{T}_{i} $.
	\end{enumerate}
\end{prop}
\begin{proof}
	Consideremos el espacio topológico producto $\displaystyle \left(\prod^{r}_{i = 1}X_{i}, \prod^{r}_{i = 1}\mathcal{T}_{i}\right) $. 
	\begin{enumerate}
	\item Si $\displaystyle V^{a_{j}}\subset X_{j} $ es entorno de $\displaystyle a_{j}\in X_{j} $, existe $\displaystyle U_{j} \in \mathcal{T}_{j} $ tal que $\displaystyle a_{j}\in U_{j}\subset V^{a_{j}} $, así tenemos que 
		\[\left(a_{1}, \ldots, a_{r}\right)\in U_{1} \times \cdots \times U_{r} \subset V^{a_{1}} \times \cdots \times V^{a_{r}} ,\]
		como $\displaystyle U_{1} \times \cdots \times U_{r} \in \prod^{r}_{i = 1}\mathcal{T}_{i} $ tenemos que $\displaystyle V^{a_{1}} \times \cdots \times V^{a_{r}} $ es entorno de $\displaystyle \left(a_{1}, \ldots, a_{r}\right) $. 
	\item Supongamos que $\displaystyle V \subset \prod^{r}_{i = 1}X_{i} $ es un entorno de $\displaystyle \left(a_{1}, \ldots , a_{r}\right) $, entonces existe $\displaystyle U\in\prod^{r}_{i = 1}\mathcal{T}_{i} $ tal que $\displaystyle \left(a_{1}, \ldots, a_{r}\right)\in U\subset V $. Como hemos visto antes, cada proyección $\displaystyle p_{j} $ es abierta por lo que $\displaystyle p_{j}\left(U\right)\in \mathcal{T}_{j} $ y $\displaystyle p_{j}\left(U\right)\subset p_{j}\left(V\right) $. 
		Así, debe ser que $\displaystyle p_{j}\left(V\right) $ es entorno de $\displaystyle a_{j}\in X_{j} $. Por tanto, existe $\displaystyle V^{a_{j}}\in \mathcal{V}\left(a_{j}\right) $ tal que $\displaystyle V^{a_{j}}\subset p_{j}\left(V\right) $. De esta forma, tendremos que por \textbf{(1)} el conjunto
		\[\left(a_{1}, \ldots, a_{r}\right)\in V^{a_{1}} \times \cdots \times V^{a_{r}} \subset V,\]
	es entorno de $\displaystyle \left(a_{1}, \ldots, a_{r}\right) $ por lo que $\displaystyle \mathcal{V}\left(a_{1}, \ldots, a_{r}\right) $ es base de entornos de $\displaystyle \left(a_{1}, \ldots, a_{r}\right) $. 	
\item Claramente todo elemento de $\displaystyle \mathcal{B}_{1 \times \cdots \times r} $ es abierto. Veamos ahora que todo abierto de la topología producto es unión de elementos de $\mathcal{B}_{1\times\cdots\times r}$. Sea
\[U\in\prod_{i=1}^r\mathcal{T}_i\]
y sea $x=(x_1,\ldots,x_r)\in U$. Por definición de la topología producto, existen abiertos $V_i\in\mathcal{T}_i$, $i=1,\ldots,r$, tales que
\[x\in V_1\times\cdots\times V_r\subset U.\]
Como $\mathcal{B}_i$ es una base de $\mathcal{T}_i$ y $x_i\in V_i$, para cada $i$ existe $U_i\in\mathcal{B}_i$ tal que
\[x_i\in U_i\subset V_i.\]
Por consiguiente,
\[x\in U_1\times\cdots\times U_r\subset V_1\times\cdots\times V_r
\subset U,\]
y además
\[U_1\times\cdots\times U_r\in\mathcal{B}_{1\times\cdots\times r}.\]
Así, para todo $x\in U$ existe un elemento de $\mathcal{B}_{1\times\cdots\times r}$ que contiene a $x$ y está contenido en $U$. Por tanto,
\[U=\bigcup_{\substack{U_1\times\cdots\times U_r\in\mathcal{B}_{1\times\cdots\times r}\\ U_1\times\cdots\times U_r\subset U}}U_1\times\cdots\times U_r.\]
Concluimos que $\mathcal{B}_{1\times\cdots\times r}$ es una base de
$\displaystyle \prod^{r}_{i = 1}\mathcal{T}_{i} $. 
	\end{enumerate}
\end{proof}

\section{Topología final}
Dado un espacio topológico $\displaystyle \left(X, \mathcal{T}\right) $ y una aplicación $\displaystyle f : X \to Y $, nos preguntamos qué topologías $\displaystyle \mathcal{T}' $ podemos considerar en $\displaystyle Y $ de modo que $\displaystyle f : \left(X, \mathcal{T}\right)\to \left(Y, \mathcal{T}'\right) $ sea continua. Una respuesta trivial sería considerar $\displaystyle \mathcal{T}' = \mathcal{T}_{\text{triv}} $. Sin embargo, buscaremos la topología más fina que cumpla la propiedad deseada.
\begin{definition}[Topología final]
Dado un espacio topológico $\displaystyle \left(X, \mathcal{T}\right) $ y una aplicación $\displaystyle f : X \to Y $, la \textbf{topología final para $\displaystyle f $} es
\[\mathcal{T}_{f} = \left\{ V \subset Y \; : \; f^{-1}\left(V\right)\in \mathcal{T}\right\}  .\]
\end{definition}
Veamos que efectivamente se trata de una topología:
\begin{itemize}
\item Tenemos que $\displaystyle f^{-1}\left(Y\right)= X \in \mathcal{T} $, por lo que $\displaystyle Y \in \mathcal{T}_{f} $ y claramente $\displaystyle \emptyset \in \mathcal{T}_{f} $. 
\item Supongamos que $\displaystyle \left\{ U_{i}\right\} _{i \in I}\subset \mathcal{T}_{f} $, entonces tendremos que 
	\[f^{-1}\left(\bigcup_{i \in I}U_{i}\right)=\bigcup_{i \in I}f^{-1}\left(U_{i}\right)\in \mathcal{T} ,\]
	por lo que $\displaystyle \bigcup_{i \in I}U_{i}\in \mathcal{T}_{f} $. 
\item Supongamos que $\displaystyle U_{1}, U_{2}\in \mathcal{T}_{f} $ no son disjuntos, entonces
	\[f^{-1}\left(U_{1} \cap U_{2}\right)=f^{-1}\left(U_{1}\right) \cap f^{-1}\left(U_{2}\right)\in \mathcal{T} .\]
	Así, debe ser que $\displaystyle U_{1} \cap U_{2} \in \mathcal{T}_{f} $. 
\end{itemize}
Veamos que es la más fina que cumple que $\displaystyle f : \left(X, \mathcal{T}\right) \to \left(Y, \mathcal{T}'\right) $ es continua. Supongamos que $\displaystyle \mathcal{T}' $ es otra topología en $\displaystyle Y $ que satisface esta propiedad. Si $\displaystyle U \in \mathcal{T}' $ tenemos que $\displaystyle f^{-1}\left(U\right)\in \mathcal{T} $ por ser $\displaystyle f $ continua, por lo que $\displaystyle U \in \mathcal{T}_{f} $. Así, $\displaystyle \mathcal{T}'\subset \mathcal{T}_{f} $.
\begin{prop}[Propiedad universal]
	Sea $\displaystyle f : \left(X, \mathcal{T}\right)\to \left(Y, \mathcal{T}'\right) $ una aplicación entre espacios topológicos. Las siguientes afirmaciones son equivalentes:
	\begin{enumerate}
	\item La topología $\displaystyle \mathcal{T}' $ en $\displaystyle Y $ es la topología final para $\displaystyle f $. 
	\item Para cada espacio topológico $\displaystyle \left(Z, \mathcal{T}''\right) $ y para cada aplicación 
		\[g : \left(Y, \mathcal{T}'\right)\to \left(Z, \mathcal{T}''\right), \]
		se tiene que $\displaystyle g $ es continua si y sólo si 
		\[g\circ f : \left(X, \mathcal{T}\right)\to \left(Z, \mathcal{T}''\right) ,\]
		es continua.
	\end{enumerate}
\end{prop}
\begin{proof}
Consideramos ambas implicaciones.
\begin{description}
\item[(1) $\displaystyle \Rightarrow $ (2)] Si $\displaystyle \mathcal{T}' = \mathcal{T}_{f} $, entonces dado $\displaystyle \left(Z, \mathcal{T}''\right) $ y $\displaystyle g : \left(Y, \mathcal{T}'\right)\to \left(Z, \mathcal{T}''\right) $,
	\begin{itemize}
	\item Si $\displaystyle g $ es continua, como $\displaystyle f $ es continua por la manera en que se ha definido la topología final, tenemos que $\displaystyle g\circ f $ es continua por ser composición de funciones continuas.
	\item Si $\displaystyle g\circ f $ es continua, entonces para cada $\displaystyle W \in \mathcal{T}'' $ tendremos que 
		\[\left(g\circ f\right)^{-1}\left(W\right)=f^{-1}\left(g^{-1}\left(W\right)\right)\in \mathcal{T} ,\]
		por lo que $\displaystyle g^{-1}\left(W\right) $ cumple la condición que caracteriza a los subconjuntos en $\displaystyle \mathcal{T}_{f} $. Por tanto, $\displaystyle g $ es continua.
	\end{itemize}
\item[(2) $\displaystyle \Rightarrow $ (1)] Tomamos $\displaystyle \left(Z, \mathcal{T}''\right) $ como $\displaystyle \left(Y, \mathcal{T}_{f}\right) $ y $\displaystyle g $ la aplicación identidad en $\displaystyle Y $, por lo que 
	\[g \circ f = id \circ f = f : \left(X, \mathcal{T}\right)\to \left(Y, \mathcal{T}_{f}\right) ,\]
	que es continua por las propiedades de la topología final para $\displaystyle f $. De esta forma,  
	\[g = id : \left(Y, \mathcal{T}'\right)\to \left(Y, \mathcal{T}_{f}\right) \]
	es continua y tenemos que $\displaystyle \mathcal{T}_{f}\subset \mathcal{T}' $. 
	Por otro lado, tomando $\displaystyle \left(Z, \mathcal{T}''\right) $ como $\displaystyle \left(Y, \mathcal{T}'\right) $ y $\displaystyle g $ como la identidad en $\displaystyle Y $,
	\[g = id : \left(Y, \mathcal{T}'\right)\to \left(Y, \mathcal{T}'\right) ,\]
	es claramente continua y por hipótesis tenemos que 
	\[g\circ f = id \circ f = f : \left(X, \mathcal{T}\right)\to \left(Y, \mathcal{T}'\right) ,\]
	debe ser continua, por lo que $\displaystyle \mathcal{T}' \subset \mathcal{T}_{f} $. 
\end{description}
\end{proof}

\begin{prop}
Un subconjunto $\displaystyle C \subset Y $ es cerrado en $\displaystyle \left(Y, \mathcal{T}_{f}\right) $ si y sólo si $\displaystyle f^{-1}\left(C\right) $ es cerrado en $\displaystyle \left(X, \mathcal{T}\right) $.
\end{prop}

\begin{proof}
En efecto, $\displaystyle C \subset Y $ es cerrado si y sólo si $\displaystyle Y\backslash C\in \mathcal{T}_{f} $ si y sólo si
\[f^{-1}\left(Y \backslash C\right)= X \backslash f^{-1}\left(C\right)\in \mathcal{T} .\]
\end{proof}
\begin{observation}
Por la proposición anterior es sencillo ver que dada $\displaystyle f : \left(X, \mathcal{T}\right)\to \left(Y, \mathcal{T}_{f}\right) $, $\displaystyle f\left(X\right)\subset Y $ es abierto y cerrado en $\displaystyle \left(Y, \mathcal{T}_{f}\right) $. 
\end{observation}

\subsection{Identificaciones}

\begin{definition}[Identificación]
Una aplicación $\displaystyle f : \left(X, \mathcal{T}\right)\to \left(Y, \mathcal{T}'\right) $ es una \textbf{identificación} si $\displaystyle f $ es sobreyectiva y $\displaystyle \mathcal{T}' = \mathcal{T}_{f} $. 
\end{definition}
\begin{observation}
Si $\displaystyle f $ es identificación entonces es continua.
\end{observation}
\begin{definition}[Conjunto saturado]
Dada una aplicación entre dos conjuntos $\displaystyle f : X \to Y $, decimos que $\displaystyle A \subset X $ es \textbf{saturado} respecto de $\displaystyle f $ o \textbf{$\displaystyle f $-saturado} si $\displaystyle f^{-1}\left(f\left(A\right)\right)=A $. 
\end{definition}

\begin{observation}
Para cada $\displaystyle B \subset Y $, el conjunto $\displaystyle A = f^{-1}\left(B\right) $ es $\displaystyle f $-saturado. En efecto, $\displaystyle f\left(A\right)=B \cap f\left(X\right) $ y $\displaystyle f^{-1}\left(B \cap f\left(X\right)\right)=f^{-1}\left(B\right) $, por lo que $\displaystyle A = f^{-1}\left(f\left(A\right)\right) $. 
\end{observation}

\begin{prop}
Sea $\displaystyle f : \left(X, \mathcal{T}\right)\to \left(Y, \mathcal{T}'\right) $ una identificación, entonces:
\begin{enumerate}
\item $\displaystyle V \subset Y $ es abierto si y sólo si $\displaystyle V = f\left(U\right) $ con $\displaystyle U \in \mathcal{T} $ y $\displaystyle U $ es $\displaystyle f $-saturado.
\item $\displaystyle C \subset Y $ es cerrado si y sólo si $\displaystyle C = f\left(H\right) $ con $\displaystyle H $ cerrado en $\displaystyle \left(X, \mathcal{T}\right) $ y $\displaystyle H $ es $\displaystyle f $-saturado.
\end{enumerate}
\end{prop}
\begin{proof}
Supongamos que $\displaystyle f : \left(X, \mathcal{T}\right)\to \left(Y, \mathcal{T}'\right) $ es una identificación.
\begin{enumerate}
\item Supongamos que $\displaystyle V \subset Y $ es abierto, entonces $\displaystyle f^{-1}\left(V\right)\in \mathcal{T}$. Podemos ver que $\displaystyle U = f^{-1}\left(V\right) $ es abierto y es un conjunto saturado. Por la sobreyectividad de $\displaystyle f $ tenemos que $\displaystyle V = f\left(U\right) $. Recíprocamente, si $\displaystyle V = f\left(U\right) $ con $\displaystyle U \in \mathcal{T} $ y $\displaystyle f $-saturado, tenemos que $\displaystyle f\left(U\right) $ es abierto puesto que $\displaystyle f^{-1}\left(f\left(U\right)\right)=U \in \mathcal{T} $. 
\item Si $\displaystyle C \subset Y $ es cerrado, entonces $\displaystyle f^{-1}\left(C\right)\subset X $ es cerrado por ser $\displaystyle f $ continua. Además, es un conjunto saturado por $\displaystyle f $. Por la sobreyectividad de $\displaystyle f $ tenemos que $\displaystyle C = f\left(f^{-1}\left(C\right)\right) $. Recíprocamente, si $\displaystyle D \subset X $ es $\displaystyle f $-saturado tenemos que $\displaystyle f\left(D\right) $ es cerrado en $\displaystyle Y $ porque $\displaystyle f^{-1}\left(f\left(D\right)\right)=D $ es cerrado en $\displaystyle X $. 
\end{enumerate}
\end{proof}
\begin{eg}
	Consideremos $\displaystyle \left(\left[0,1\right] , \mathcal{T}|_{\left[0,1\right] }\right) $ y la aplicación 
	\[f : \left[0,1\right] \to \mathbb{S}^{1} : t \to \left(\cos2\pi t, \sin2\pi t\right) .\]
	Si dotamos a $\displaystyle \mathbb{S}^{1} $ con la topología final para $\displaystyle f $ obtenemos la topología que tiene como subespacio de $\displaystyle \left(\R^{2}, \mathcal{T}_{u}\right) $. \\

	En efecto, podemos observar que si usamos la topología usual en $\displaystyle \R^{2} $ $\displaystyle f $ es continua, por lo que debe ser que $\displaystyle \mathcal{T}_{u} | _{\mathbb{S}^{1}}\subset \mathcal{T}_{f} $. Ahora, supongamos que $\displaystyle A \in \mathcal{T}_{f} $. Por definición tenemos que $\displaystyle f^{-1}\left(A\right)\subset\left[0,1\right]  $ es abierto en $\displaystyle \mathcal{T}_{u}|_{\left[0,1\right] } $. Podemos distinguir dos casos:
	\begin{itemize}
	\item Si $\displaystyle \left(1,0\right)\not\in A $ tenemos que $\displaystyle 0,1\not\in f^{-1}\left(A\right) $. En este caso tenemos que $\displaystyle f^{-1}\left(A\right) $ es abierto en $\displaystyle \left(0,1\right) $ por lo que es abierto en $\displaystyle \R $. Así, se puede expresar como unión de intervalos abiertos $\displaystyle \left(a,b\right)\subset \left(0,1\right) $ para cada uno de los cuales su imagen es un abierto
		\[f\left(a,b\right)=\Omega_{\left(a,b\right)}\cap \mathbb{S}^{1}\in \mathcal{T}_{u}|_{\mathbb{S}^{1}} ,\]
		con $\displaystyle \Omega_{\left(a,b\right)} $ la región abierta del plano usual comprendida entre las rectas $\displaystyle r_{a} $ y $\displaystyle r_{b} $ que unen, respectivamente, el origen de coordenadas con los puntos $\displaystyle f\left(a\right) $ y $\displaystyle f\left(b\right) $. Por tanto, $\displaystyle A \in \mathcal{T}_{u}|_{\mathbb{S}^{1}} $.
	\item Si $\displaystyle \left(1,0\right)\in A $ tenemos que $\displaystyle 0,1\in f^{-1}\left(A\right) $. Así, existe $\displaystyle 0 < \epsilon < \frac{1}{4} $ tal que $\displaystyle [0,\epsilon) \cup(1-\epsilon, 1 ] \subset f^{-1}\left(A\right) $. Por tanto, $\displaystyle \left(1,0\right)\in f\left([0,\epsilon) \cup(1-\epsilon,1]\right)\subset A $, pero 
		\[f\left([0,\epsilon)\cup(1-\epsilon,1]\right)= \left\{ \tan\left(-2\pi \epsilon \right)x < y < \tan 2\pi \epsilon x, \; x > 0\right\} \cap \mathbb{S}^{1} \in \mathcal{T}_{u} |_{\mathbb{S}^{1}} .\]
		Como $\displaystyle f^{-1}\left(A\right)\backslash \left\{ 0,1\right\}  $ es abierto usual en $\displaystyle \left(0,1\right) $, a la vista del caso anterior concluimos que $\displaystyle A $ es abierto en $\displaystyle \mathbb{S}^{1} $. 
	\end{itemize}
\end{eg}
\begin{eg}
Consideramos el subconjunto 
\[ C = \left\{ x^{2}+y^{2}=1, \; 0 \leq z \leq 1\right\} \subset \R^{3} ,\]
y la aplicación 
\[f : \left[0,1\right] ^{2} \to C : \left(s,t\right)\to \left(\cos 2\pi s, \sin 2\pi s, t\right) .\]
Tomando la topología final para $\displaystyle f $ en $\displaystyle C $, obtenemos la misma topología que la que hereda $\displaystyle C $ de la usual en $\displaystyle \R^{3} $. Esto se puede razonar de modo análogo a como lo hicimos en el ejemplo anterior. Tampoco en este caso la aplicación $\displaystyle f $ es abierta. 
\end{eg}
\begin{prop}
	Sea $\displaystyle f : \left(X, \mathcal{T}\right)\to \left(Y, \mathcal{T}'\right) $ una aplicación. 
	\begin{enumerate}
	\item Si $\displaystyle f $ es sobreyectiva, continua y abierta, es una identificación.
	\item Si $\displaystyle f $ es sobreyectiva, continua y cerrada, es una identificación.
	\end{enumerate}
\end{prop}
\begin{proof}
Supongamos que $\displaystyle f : \left(X, \mathcal{T}\right)\to \left(Y, \mathcal{T}'\right) $ es una aplicación. 
\begin{enumerate}
\item Sólamente tenemos que ver que $\displaystyle \mathcal{T}'=\mathcal{T}_{f} $. Dado $\displaystyle V \subset Y $ tal que $\displaystyle f^{-1}\left(V\right)\subset X $ es abierto, tendremos que por ser $\displaystyle f $ abierta, $\displaystyle f\left(f^{-1}\left(V\right)\right)\in \mathcal{T}' $. Por la sobreyectividad de $\displaystyle f $, $\displaystyle V = f\left(f^{-1}\left(V\right)\right) $. Por tanto, $\displaystyle \mathcal{T}_{f}\subset \mathcal{T}' $. La continuidad de $\displaystyle f $ nos garantiza el otro contenido, por lo que $\displaystyle \mathcal{T}' = \mathcal{T}_{f} $. 
\item Dado $\displaystyle V \subset Y $ para el que $\displaystyle f^{-1}\left(V\right)\subset X $ sea un abierto, se tiene por tanto que $\displaystyle X\backslash f^{-1}\left(V\right) $ es cerrado en $\displaystyle X $. Como $\displaystyle f $ es cerrada, $\displaystyle f\left(X\backslash f^{-1}\left(V\right)\right) $ es un cerrado en $\displaystyle Y $, pero como $\displaystyle f $ es sobreyectiva $\displaystyle f\left(X \backslash f^{-1}\left(V\right)\right)=Y \backslash V $. Por tanto $\displaystyle V \in \mathcal{T}' $ por lo que $\displaystyle \mathcal{T}_{f}\subset \mathcal{T}' $. El otro contenido viene dado por la continuidad de $\displaystyle f $. 
\end{enumerate}
\end{proof}

\begin{observation}
Hay identificaciones que son abiertas y cerradas; identificaciones abiertas pero no cerradas; identificaciones no abiertas pero sí cerradas; identificaciones no abiertas y no cerradas.
\end{observation}
\subsection{Cocientes topológicos}
Supongamos que tenemos un espacio topológico $\displaystyle \left(X, \mathcal{T}\right) $ y que en $\displaystyle X $ tenemos definida una relación de equivalencia $\displaystyle \sim  $. Sea $\displaystyle Y = X / \sim  $, entonces podemos definir una aplicación natural
\[p : X \to Y : x \to \left[x\right]  .\]
\begin{definition}[Espacio topológico cociente]
Si consideramos en el conjunto cociente $\displaystyle Y $ la topología final para la aplicación $\displaystyle p $, obtenemos el \textbf{espacio topológico cociente} $\displaystyle \left(Y, \mathcal{T}_{p}\right) $. Decimos que $\displaystyle \mathcal{T}_{p} $ es la \textbf{topología cociente} de la de $\displaystyle X $.
\end{definition}

Hay una correspondencia entre identificaciones y cocientes topológicos. Sea $\displaystyle f : \left(X, \mathcal{T}\right)\to \left(Y, \mathcal{T}'\right) $ una aplicación sobreyectiva entre dos espacios topológicos. Definimos en $\displaystyle X $ la relación de equivalentcia
\[x_{1} \sim x_{2} \iff f\left(x_{1}\right)=f\left(x_{2}\right) .\]
Podemos obtener la factorización canónica de $\displaystyle f $ como la composición de la aplicación $\displaystyle p : X \to X / \sim  $ y de 
\[\overline{f} : X / \sim \; \to Y : \left[x\right] \to f\left(x\right) .\]
Claramente esta función está bien definida y es biyectiva.
\begin{prop}
En la situación que acabamos de describir, considerando en $\displaystyle X/\sim $ la topología cociente, se tiene que $\displaystyle f $ es una identificación si y sólo si $\displaystyle \overline{f} $ es homeomorfismo.
\end{prop}
\begin{proof}
Empezamos observando que $\displaystyle A \subset X $ es saturado por $\displaystyle f $ si y sólo si es saturado por $\displaystyle p $. En efecto, la relación $\displaystyle \sim $ nos da
\[p^{-1}\left(p\left(A\right)\right)= \left\{ x \in X \; : \; \exists a \in A, \; x \sim a\right\} = \left\{ x \in X \; : \; \exists a \in A, \; f\left(x\right)=f\left(a\right),\right\} = f^{-1}\left(f\left(A\right)\right) .\]
\begin{description}
\item[($\displaystyle \Rightarrow $)] Supongamos que $\displaystyle f $ es una identificación, entonces dado $\displaystyle V \subset Y $ abierto tendremos que 
\[\overline{f}^{-1}\left(V\right)=p\left(f^{-1}\left(V\right)\right) ,\]
que es abierto en $\displaystyle X/\sim $, porque es la imagen por $\displaystyle p $ del abierto $\displaystyle f^{-1}\left(V\right) $ que es saturado por $\displaystyle p $. Concluimos que $\displaystyle \overline{f} $ es continua. 
Por otra parte, $\displaystyle \overline{f} $ es abierta. En efecto, un abierto en $\displaystyle X / \sim  $ es la imagen por $\displaystyle p $ de un abierto $\displaystyle U \subset X $ saturado por $\displaystyle p $. Entonces, $\displaystyle \overline{f}\left(p\left(U\right)\right)=f\left(U\right) $, que es abierto en $\displaystyle Y $ ya que $\displaystyle f^{-1}\left(f\left(U\right)\right)=U $.
\item[($\displaystyle \Leftarrow $)] Supongamos ahora que $\displaystyle \overline{f} $ es homeomorfismo. Claramente $\displaystyle f $ es sobreyectiva. Por la propiedad universal de la topología final (aplicada a la topología cociente), $\displaystyle f $ es continua, por lo que $\displaystyle \mathcal{T}' \subset \mathcal{T}_{f} $. Dado $\displaystyle B \subset Y $ con $\displaystyle f^{-1}\left(B\right)\in \mathcal{T} $, tenemos que $\displaystyle p\left(f^{-1}\left(B\right)\right) $ es abierto pues $\displaystyle f^{-1}\left(B\right) $ es saturado por $\displaystyle p $. Por tanto $\displaystyle B = \overline{f}\left(p\left(f^{-1}\left(B\right)\right)\right) $ es abierto en $\displaystyle Y $.
\end{description}
\end{proof}
\begin{observation}
En la situación que acabamos de ver, $\displaystyle Y $ se puede interpretar como un modelo del espacio cociente $\displaystyle X/\sim  $. En general, los cocientes son muy útiles para definir nuevos espacios topológicos. Disponer de un modelo y de una identificación es muy conveniente para estudiar las propiedades de esos cocientes.
\end{observation}
\begin{eg}
	La identificación $\displaystyle f : \left[0,1\right] \to \mathbb{S}^{1} $ nos da un modelo para el cociente $\displaystyle \left[0,1\right] / \sim  $ donde $\displaystyle 0 \sim 1 $ y para $\displaystyle t \in \left(0,1\right) $, $\displaystyle t $ está relacionado con él mismo.
\end{eg}
\begin{eg}
	La identificación $\displaystyle f : \left[0,1\right] \times \left[0,1\right] \to C $ nos da un modelo para el cociente $\displaystyle \left[0,1\right] \times\left[0,1\right] /\sim $, donde $\displaystyle \left(0,t\right)\sim\left(1,t\right) $ para cada $\displaystyle t \in \left[0,1\right]  $, mientras que para cualquier $\displaystyle s \in \left(0,1\right) $ el punto $\displaystyle \left(s,t\right) $ sólo está relacionado con él mismo.
\end{eg}
\subsection{Sumas finitas}
Dado un número finito de espacios topológicos $\displaystyle \left\{ \left(X_{i}, \mathcal{T}_{i}\right)\right\} _{i = 1, \ldots, r} $, consideramos el conjunto
\[\sum^{r}_{i= 1}X_{i} = X_{1} + \cdots + X_{r} : = \left(X_{1} \times \left\{ 1\right\} \right) \cup \cdots \cup \left(X_{r} \times r\right) .\]
Podemos pensar en este conjunto como la unión disjunta de los $\displaystyle X_{i} $. Para cada $\displaystyle i \in \left\{ 1, \ldots, r\right\}  $ podemos considerar la aplicación, 
\[e_{i} : \left(X_{i}, \mathcal{T}_{i}\right)\to \sum^{r}_{j = 1}X_{j} : x_{i} \to \left(x_{i}, i\right) .\]
Queremos dotar a $\displaystyle \sum^{r}_{j= 1}X_{j} $ de la topología más fina que haga continuas todas las aplicaciones $\displaystyle e_{i} $. La solución a este problema es la que se define a continuación:
\begin{definition}[Topología suma]
	El \textbf{espacio topológico suma} de los espacios $\displaystyle \left(X_{i}, \mathcal{T}_{i}\right)  $ para $\displaystyle i = \left\{ 1, \ldots, r\right\}  $ está formado por el conjunto $\displaystyle \sum^{r}_{i = 1}X_{i} $ en el que consideramos la \textbf{topología suma} de las topologías $\displaystyle \mathcal{T}_{i} $, que es la que tiene por base la siguiente familia de subconjuntos
	\[ \mathcal{B} = \left\{ U_{i} \times \left\{ i\right\} \; : \; U_{i} \in \mathcal{T}_{i}, i \in \left\{ 1, \ldots, r\right\} \right\}  .\]
	Denotaremos dicha topología por $\displaystyle \sum^{r}_{i=1}\mathcal{T}_{i} $. 
\end{definition}
Veamos que esta es la base de una topología en $\displaystyle \sum^{r}_{i=1}X_{i} $ y además es la más fina que consigue que todas las $\displaystyle e_{i} $ sean continuas. 
\begin{itemize}
	\item Claramente tenemos que $\displaystyle X_{i} \times \left\{ i\right\} \in \mathcal{B} $, $\displaystyle \forall i \in \left\{ 1, \ldots, r\right\}  $, por lo que 
		\[\sum^{r}_{i = 1}X_{i} = (X_{1} \times \left\{ 1\right\}) \cup \cdots \cup (X_{r} \times \left\{ r\right\}) \in \mathcal{T}_{\mathcal{B}} .\]
	\item Supongamos ahora que $\displaystyle B_{1}, B_{2} \in \mathcal{B} $ y $\displaystyle x \in B_{1} \cap B_{2} $. Si $\displaystyle B_{i} = \left(U_{1}^{i} \times \left\{ 1\right\} \right)\cup \cdots \cup \left(U_{r}^{i}\times \left\{ r\right\} \right)  $ para $\displaystyle i \in \left\{ 1,2\right\}  $, tenemos que 
		\[B_{1} \cap B_{2} = \left(\left(U_{1}^{1}\cap U_{1}^{2}\right)\times \left\{ 1\right\} \right) \cup \cdots \cup \left(\left(U_{r}^{1}\cap U_{r}^{2}\right)\times \left\{ r\right\} \right) \in \mathcal{B}.\]
\end{itemize}
Por tanto, tenemos que $\displaystyle \mathcal{B} $ es una base de una topología en $\displaystyle \mathcal{T}_{\mathcal{B}} $. Es sencillo ver que $\displaystyle \mathcal{T}_{\mathcal{B}} $ hace que todas las aplicaciones $\displaystyle e_{i} $ sean continuas. Ahora, veamos que es la topología más fina que hace continuas a $\displaystyle e_{i} $. Supongamos que $\displaystyle \mathcal{T}' $ también cumple esta propiedad. Si $\displaystyle U \in \mathcal{T}' $ tenemos que $\displaystyle e_{i}^{-1}\left(U\right)\in \mathcal{T}_{i} $, por lo que 
\[e^{-1}_{i} \left(U\right)=\bigcup_{j \in J}V_{j}, \quad V_{j} \in \mathcal{T}_{i}.\]
Como cada $\displaystyle e_{i} $ es inyectiva, $\displaystyle e_{i}\left(e_{i}^{-1}\left(U\right)\right)= \left(\bigcup_{j \in J_{i}}V_{j}\right) \times \left\{ i\right\} $, de forma que 
\[U = \bigsqcup_{1 \leq i \leq r}e_{i}\left(e^{-1}_{i}\left(U\right)\right) =\bigsqcup_{1 \leq i \leq r} \left(\left(\bigcup_{j \in J_{i}}V_{j}\right) \times \left\{ i\right\}\right) \in \mathcal{T}_{\mathcal{B}}.\]
De esta forma, $\displaystyle \mathcal{T}'\subset\mathcal{T}_{\mathcal{B}} $. 
\begin{prop}
	La aplicación 
	\[e_{i} : \left(X_{i}, \mathcal{T}_{i}\right)\to \left(\sum^{r}_{j = 1}X_{j}, \sum^{r}_{j = 1}\mathcal{T}_{j}\right) \]
	es un embebimiento topológico, es abierta y es cerrada para cada $\displaystyle i \in \left\{ 1, \ldots, r\right\}  $. 
\end{prop}
\begin{proof}
Basta con demostrar que 
\[e_{i} : \left(X_{i}, \mathcal{T}_{i}\right)\to \left(e_{i}\left(X_{i}\right), \left(\sum^{r}_{j = 1}\mathcal{T}_{j}\right)|_{e_{i}\left(X_{i}\right)}\right) ,\]
es un homeomorfismo. Está claro que es biyectiva. Además es abierta, puesto que para cada $\displaystyle U_{i} \in \mathcal{T}_{i} $ tenemos que $\displaystyle e_{i}\left(U_{i}\right)=U_{i} \times \left\{ i\right\} \in \sum^{r}_{j = 1}\mathcal{T}_{j} $ y está contenido en $\displaystyle e_{i}\left(X_{i}\right)=X_{i} \times \left\{ i\right\}  $. Está claro que por definición de topología suma $\displaystyle e_{i} $ es continua. Finalmente, para ver que es cerrada podemos ver que si $\displaystyle C_{i} \subset X_{i} $ es cerrado, entonces
\[\left(\sum^{r}_{j=1}X_{i}\right)\backslash e_{i}\left(C_{i}\right)=\bigcup_{j \neq i}\left(X_{j} \times \left\{ j\right\} \right) \cup \left(\left(X_{i}\backslash C_{i}\right)\times \left\{ i\right\} \right)\in \sum^{r}_{j = 1}\mathcal{T}_{j} .\]
\end{proof}
\begin{observation}
Podemos ver que, en particular, los espacios $\displaystyle X_{i} $ son subespacios abiertos y cerrados de la suma. 
\end{observation}
\begin{prop}[Propiedad universal]
	Dada una familia finita de espacios topológicos $\displaystyle \left\{ \left(X_{i}, \mathcal{T}_{i}\right)\right\} _{i \in \left\{ 1, \ldots, r\right\} } $, dotemos a $\displaystyle \sum^{r}_{i=1}X_{i} $ de una topología $\displaystyle \mathcal{T}' $. Entonces, las siguientes afirmaciones son equivalentes:
	\begin{enumerate}
	\item $\displaystyle \mathcal{T}' = \sum^{r}_{i = 1}\mathcal{T}_{i} $.
	\item Para cada espacio topológico $\displaystyle \left(Z, \mathcal{T}''\right) $ y para cada aplicación
		\[g : \left(\sum^{r}_{i = 1}X_{i}, \mathcal{T}'\right)\to \left(Z, \mathcal{T}''\right) ,\]
		$\displaystyle g $ es continua si y sólo si $\displaystyle g \circ e_{i} : \left(X_{i}, \mathcal{T}_{i}\right)\to \left(Z, \mathcal{T}''\right)$ es continua, para $\displaystyle i = 1, \ldots, r $. 
	\end{enumerate}
\end{prop}
\begin{proof}
	Demostramos ambas implicaciones. 
	\begin{description}
	\item[($\displaystyle \Rightarrow $)] Supongamos que $\displaystyle \mathcal{T}' = \sum^{r}_{i = 1}\mathcal{T}_{i} $. 
		\begin{itemize}
		\item Si $\displaystyle g $ es continua, $\displaystyle g\circ e_{i} $ es continua por ser composición de funciones continuas. 
		\item Si $\displaystyle g\circ e_{i} $ es contiua, para $\displaystyle U \in \mathcal{T}'' $ tenemos que 
			\[\left(g\circ e_{i}\right)^{-1}\left(U\right)=e^{-1}_{i}\left(g^{-1}\left(U\right)\right)\in \mathcal{T}_{i} ,\]
		
	\end{itemize}
		
	\end{description}
\end{proof}

Como idea general, el espacio $\displaystyle \sum^{r}_{j = 1}X_{j} $ es, localmente, como cada uno de los espacios (es decir, $\displaystyle \sum^{r}_{j = 1}X_{j} $ es unión disjunta de subespacios) que son cada uno de ellos como uno de los espacios de partida. Así, para determinar bases de entornos en $\displaystyle \sum^{r}_{j = 1}X_{j} $, consideramos para cada uno de sus puntos una base de entornos en el espacio $\displaystyle X_{i} $ al que pertenece. \\

Globalmente, se trata cada sumando por separado. De esta forma, una base de la topología del espacio suma se puede obtener uniendo bases de abiertos de cada uno de los espacios sumando. 
\subsection{Espacios proyectivos reales}
Recordamos que podemos definir el espacio proyectivo real de dimensión $\displaystyle n $, que denotamos por $\displaystyle \mathbb{P}^{n}_{\R} $, como el conjunto cociente $\displaystyle \mathbb{P}^{n}_{\R} = \left(\R^{n+1}\backslash \left\{ 0\right\} \right)/ \sim $, donde $\displaystyle \sim  $ es la relación de equivalencia definida por
\[\left(x_{0}, \ldots, x_{n}\right)\sim \left(y_{0}, \ldots, y_{n}\right)\iff \exists\lambda\neq 0, \; \left(y_{0}, \ldots, y_{n}\right)=\lambda\left(x_{0}, \ldots, x_{n}\right) .\]
Consideramos el espacio proyectivo real como espacio topológico, dotándole de la topología cociente, esto es, la topología final para la aplicación
\[\pi : \R^{n+1} \backslash \left\{ 0\right\} \to \mathbb{P}^{n}_{\R} : \left(x_{0}, \ldots, x_{n}\right) \to \left[x_{0} : \cdots : x_{n}\right]  .\]
Podemos ver los puntos del espacio proyectivo como rayos, es decir, subespacios vectoriales de dimensión uno en $\displaystyle \R^{n+1} $. \\
Consideremos un hiperplano proyectivo $\displaystyle H \subset \mathbb{P}^{n}_{\R} $ que viene determinado por un hiperplano vectorial $\displaystyle \hat{H}\subset \R^{n+1} $, que corresponde al subespacio en el que se anula una forma lineal $\displaystyle h $. En otras palabras, $\displaystyle H = \left(\hat{H} \backslash \left\{ 0\right\} \right)/ \sim $ con $\displaystyle \hat{H}=\Ker h $. Observamos que $\displaystyle H $ es cerrado en $\displaystyle \mathbb{P}^{n}_{\R} $. Su complementario $\displaystyle U = \mathbb{P}^{n}_{\R} \backslash  H $ es un abierto, de hecho
\[U = \pi\left(\R^{n+1} \backslash \hat{H}\right), \quad \hat{H} = \left\{ h = 0\right\} \subset \R^{n+1} ,\]
y $\displaystyle \R^{n+1} \backslash \hat{H} $ es un abierto $\displaystyle \pi  $-saturado en $\displaystyle \R^{n+1} \backslash \left\{ 0\right\}  $. 
\begin{definition}[Carta afín]
Cada subconjunto $\displaystyle U $ definido como acabamos de ver se llama una \textbf{carta afín} de $\displaystyle \mathbb{P}^{n}_{\R} $. El hiperplano $\displaystyle H $ se dice que es el \textbf{hiperplano de infinito} de la carta afín $\displaystyle U $.
\end{definition}
Si ahora nos fijamos en el hiperplano afín $\displaystyle \left\{ h = 1\right\} \subset \R^{n+1} \backslash \left\{ 0\right\}  $, podemos establecer una biyección restringiendo la aplicación $\displaystyle \pi  $,
\[\pi |_{ \left\{ h = 1\right\} } : \left\{ h = 1\right\} \to \mathbb{P}^{n}_{\R} \backslash H = U .\]
\begin{prop}
	La aplicación $\displaystyle \pi  $ restringida antes mencionada, considerando en $\displaystyle \left\{ h =1\right\}  $ la topología relativa como subespacio de $\displaystyle \R^{n+1} $ y en $\displaystyle U $ la topología relativa como subespacio de $\displaystyle \mathbb{P}^{n}_{\R} $, es un homeomorfismo.
\end{prop}
\begin{proof}
	Sabemos que $\displaystyle \pi |_{ \left\{ h = 1\right\} } $ es biyectiva. Veamos que es continua. Como $\displaystyle U \subset \mathbb{P}^{n}_{\R} $ es abierto, todo abierto relativo $\displaystyle V \subset U $ es abierto en $\displaystyle \mathbb{P}^{n}_{\R} $. Por la definición de la topología cociente, tenemos entonces que $\displaystyle \pi ^{-1}\left(V\right)\subset \R^{n+1} \backslash \left\{ 0\right\}  $ es abierto. Por tanto, 
	\[ \left(\pi|_{ \left\{ h=1\right\} }\right)^{-1}\left(V\right) =\pi^{-1}\left(V\right)\cap \left\{ h = 1\right\}  ,\]
	es abierto en $\displaystyle \left\{ h=1\right\}  $ como subespacio de $\displaystyle \R^{n+1} $. Por otra parte, $\displaystyle \pi| _{ \left\{ h = 1\right\} } $ es abierta. En efecto, si tomamos un abierto básico de $\displaystyle \left\{ h = 1\right\}  $ como subespacio de $\displaystyle \R^{n+1} $, un disco abierto $\displaystyle D $, su imagen es el conjunto de subespacios de dimensión uno determinados por todos los vectores de $\displaystyle \R^{n+1} $ correspondientes a los puntos de $\displaystyle D $, es decir, los rayos del cono sobre $\displaystyle D $. Este conjunto es abierto en $\displaystyle U $ y en $\displaystyle \mathbb{P}^{n}_{\R} $. 
\end{proof}

Podemos pensar en $\displaystyle \mathbb{P}^{n}_{\R} $ viéndolo a partir del espacio suma de cartas afines. En concreto, basta tomar $\displaystyle n + 1 $ cartas distintas que cubren todo $\displaystyle \mathbb{P}^{n}_{\R} $. Habitualmente, consideramos las cartas afines
\[U_{i}= \left\{ \left[x_{0} : \cdots : x_{n}\right] \in \mathbb{P}^{n}_{\R} \; : \; x_{i} \neq 0\right\} , \; i \in \left\{ 0, \ldots, n\right\}  .\]
Para entender el espacio proyectivo tenemos que pegar esas cartas afines. Para ello es fundamental la siguiente observación. 
\begin{observation}
	Si $\displaystyle U_{1} $ y $\displaystyle U_{2} $ son dos cartas afines distintas, con hiperplanos de infinito $\displaystyle \left\{ h_{1}=0\right\}  $ y $\displaystyle \left\{ h_{2}=0\right\}  $, respectivamente, tenemos que $\displaystyle U_{1} \cap U_{2} $ es abierto en $\displaystyle U_{1} $ y en $\displaystyle U_{2} $. La topología heredada como subespacio es la misma en ambos casos. \\

	Veámoslo usando los homeomorfismos con espacios afines. Tenemos que $\displaystyle U_{1} $ es homeomorfo a $\displaystyle \left\{ h_{1} = 1\right\}  $. Análogamente, $\displaystyle U_{2}  $ es homeomorfo a $\displaystyle \left\{ h_{2}=1\right\}  $. En $\displaystyle \left\{ h_{1}=1\right\}  $, el abierto correspondiente a $\displaystyle U_{1} \cap U_{2} $ está definido por $\displaystyle h_{2}\neq 0 $. En $\displaystyle \left\{ h_{2}=1\right\}  $ el abierto correspondiente a $\displaystyle U_{1}\cap U_{2} $ está definido por $\displaystyle h_{1}\neq 0 $. Podemos establecer un homeomorfismo explícito entre estos dos abiertos. \\

	Veámoslo para simplificar, tomando $\displaystyle U_{1}= \left\{ x_{1}\neq 0\right\}  $ y $\displaystyle U_{2} = \left\{ x_{2}\neq 0\right\}  $ en $\displaystyle \mathbb{P}^{n}_{\R} $. Entonces, las coordenadas en $\displaystyle U_{1} $ visto como hiperplano afín $\displaystyle \left\{ x_{1} = 1\right\}  $ son $\displaystyle \left(\frac{x_{0}}{x_{1}}, 1, \frac{x_{2}}{x_{1}}, \ldots, \frac{x_{n}}{x_{1}}\right) $, que podemos renombrar $\displaystyle \left(u_{1}, 1, \ldots, u_{n}\right) $. Análogamente, en $\displaystyle U_{2} $ visto como hiperplano afín $\displaystyle \left\{ x_{2}=1\right\}  $ son
	$\displaystyle \left(\frac{x_{0}}{x_{2}}, \frac{x_{1}}{x_{2}}, 1, \ldots, \frac{x_{n}}{x_{2}}\right) $, que podemos renombrar $\displaystyle \left(\psi_{1}, \psi_{2}, 1, \ldots, \psi_{n}\right) $. El homeomorfismo buscado será
	\[\varphi_{12} : \left\{ x_{1} = 1, \; x_{2}\neq 0\right\} \to \left\{ x_{1}\neq 0, \; x_{2}=1\right\} : \left(u_{1}, 1, \ldots, u_{n}\right)\to \left(\frac{u_{1}}{u_{2}}, \frac{1}{u_{2}}, 1, \ldots, \frac{u_{n}}{u_{2}}\right) .\]
\end{observation}

